\section{Related Work}
\label{sec:related_work}
\textbf{LLM-authored weights.} \citet{Takehi2026SparksWeightspace} coined weight-space programming and report gains on a 0.8B model. The closest prior work is \citet{Bennington2026GradientfreeWeight}. They cut thinking tokens in Qwen3.8-27B with rank-one edits and test against a norm-matched orthogonal random control, which we adopt (TACT~\citep{Sui2026TactMitigating} also reports one). Both evaluate their edits on single-turn benchmarks. \citet{Mehta2026InterpretableWeightspace} have an LLM agent write the weights of small transformers and brain encoding models from scratch. The agent does no gradient training; a fixed evaluator fits the encoding models' linear read-out. They audit each saved model for interpretability. Their models are built whole for a fixed evaluator. We make sparse closed-form edits to a pretrained 4B-parameter agent, each with a written prediction and a matched random control, and evaluate whole agent episodes. We keep the idea of Takehi et al.\ and Bennington et al.\ and the latter's control; \S\ref{sec:intro} lists what we change.

% Comparison with closest prior work. Every cell is sourced in related_table_sources.md.
\begin{table}[t]
\centering
\small
\caption{How this work differs from the closest prior work. \emph{Cond.}: the change depends on a condition read from the hidden state. \emph{Plain weights}: the change is an ordinary weight tensor, servable by a stock engine without hooks. \emph{Control}: a random-direction or random-vector control reported. \emph{Agents}: evaluated on multi-step agent episodes. \emph{Author}: who produces the change (``contrast'': a direction from activation contrasts on human-chosen data, with layer and strength set by the authors). $^a$Random vectors of similar magnitude, qualitative and KL comparison. $^b$TACT-Gate and TACT-Cap; TACT-Add is unconditional. $^c$The rank-one update responds to the subject's key vector. $^d$For the three unconditional edits (\S\ref{sec:control}--\ref{sec:devgain}); in the benchmark each author built its own random control (Claude's is ungated and smaller than its gated edits; GPT's omits the rank-two rotation). $^e$For brain encoding the agent writes the feature network and a fixed evaluator fits a ridge read-out. $^\dagger$Dense weights edited in place are plain weights by construction; the source does not report serving them.}
\label{tab:related}
\resizebox{\linewidth}{!}{%
\begin{tabular}{@{}lllllll@{}}
\toprule
Work & Lives in & Cond. & Plain weights & Control & Agents & Author \\
\midrule
ActAdd~\citep{Turner2023SteeringLanguage} & Hook & No & No & Random vector$^a$ & No & Human (contrast) \\
CAA~\citep{Panickssery2023SteeringLlama} & Hook & No & No & No & No & Human (contrast) \\
CAST~\citep{Lee2024ProgrammingRefusal} & Hook & Yes & No & No & No & Human (contrast) + grid search \\
TACT~\citep{Sui2026TactMitigating} & Hook & Yes$^b$ & No & Norm-matched random & Yes & Human (contrast) \\
ROME~\citep{Meng2022LocatingEditing} & Weights & Key$^c$ & Yes$^\dagger$ & No & No & Optimisation + closed form \\
\citet{Arditi2024RefusalLanguage} & Weights & No & Yes$^\dagger$ & No & No & Human (contrast) \\
\citet{Takehi2026SparksWeightspace} & Weights & No & Yes$^\dagger$ & No & No & LLM author \\
\citet{Bennington2026GradientfreeWeight} & Weights & No & Yes & Norm-matched random & No & Human (contrast) \\
\citet{Mehta2026InterpretableWeightspace} & Weights (whole small model)$^e$ & n/a & n/a & Interpretable baselines & No & LLM author \\
\midrule
This paper & Weights (LoRA) & Yes (gate) & Yes (vLLM) & Norm-matched random$^d$ & Yes & LLM author \\
\bottomrule
\end{tabular}%
}
\end{table}


\textbf{Activation steering.} Steering methods add a direction to the residual stream at inference. Activation addition~\citep{Turner2023SteeringLanguage}, contrastive activation addition~\citep{Panickssery2023SteeringLlama} and representation engineering~\citep{Zou2023RepresentationEngineering} obtain the direction from contrasting prompts, as we do. Inference-time intervention~\citep{Li2023InferencetimeIntervention} and in-context vectors~\citep{Liu2023IncontextVectors} are close variants. Conditional activation steering~\citep{Lee2024ProgrammingRefusal} applies a steering vector only when a condition direction is present in the hidden state. It is the inference-time counterpart of our gate. ReFT~\citep{Wu2024ReftRepresentation} learns such interventions by training. AxBench~\citep{Wu2025AxbenchSteering} finds that prompting beats representation-based steering and that difference-of-means is the strongest concept detector. TACT~\citep{Sui2026TactMitigating} steers coding agents along overthinking and overacting axes and reports gains of 4.8--5.8 points. \citet{Silva2026LatentProgramming} show that linear probes on a coding agent's hidden states forecast whether its pending edits will succeed. Our probe for patch correctness (\S\ref{sec:negative}) is a signal of the same kind. All of these methods need hooks or extra modules in the inference stack. We compile the direction and its condition into existing weights, because the competition scores submitted weights served by a stock engine (\S\ref{sec:primitive}).

\textbf{Direct weight editing.} ROME~\citep{Meng2022LocatingEditing} and MEMIT~\citep{Meng2022MasseditingMemory} change factual associations with closed-form updates to MLP weights. \citet{Arditi2024RefusalLanguage} remove a refusal direction by orthogonalising weights against it. These edits remove or replace what existing weights write. Our primitive adds a new write by repurposing a neuron that contributes almost nothing, and it targets a recurring agent decision.

\textbf{Small trained changes and agent training.} Deletion Tuning~\citep{Lu2026SurprisingEffectiveness}, BitFit~\citep{Benzaken2021BitfitSimple} and task arithmetic~\citep{Ilharco2022EditingModels} change few weights but obtain the change by training. Our edits and our one-shot deletion experiment (\S\ref{sec:negative}) use no optimisation step. For software agents, SWE-bench~\citep{Jimenez2023SwebenchCan} defines the task, and SWE-agent~\citep{Yang2024SweagentAgentcomputer} defines an agent--computer interface for it. SWE-Gym~\citep{Pan2024TrainingSoftware}, R2E-Gym~\citep{Jain2025R2egymProcedural} and SWE-smith~\citep{Yang2025SwesmithScaling} improve agents by training on executable repository environments. SWE-Gym uses real environments; R2E-Gym and SWE-smith synthesise them.
